\documentclass[11pt,letterpaper]{article}
% =====================================================================
%  FangMath -- FLASH CARDS
%  Linear Functions: Parallel and Perpendicular Lines (2 cards)
%  Page 1 = fronts, Page 2 = backs (same top/bottom order).
%  Print double-sided, "flip on long edge", then cut along the dashed line.
%  Single-sided fallback: print both pages, cut, and glue back to back.
% =====================================================================
\usepackage[T1]{fontenc}
\usepackage[scaled]{helvet}
\renewcommand{\familydefault}{\sfdefault}
\usepackage{amsmath,amssymb}
\usepackage[margin=0.6in,top=0.55in,bottom=0.7in]{geometry}
\usepackage[most]{tcolorbox}
\usepackage{tikz}
\usepackage{array}
\usepackage{fancyhdr}

% ---- Palette "F. Bright & Playful" ----
\definecolor{fmtitle}{HTML}{3D2C8D}   % titles / card fronts
\definecolor{fmkey}{HTML}{0096C7}     % Key Idea / card backs
\definecolor{fmwatch}{HTML}{2BB673}   % Watch Me
\definecolor{fmturn}{HTML}{F4A100}    % Your Turn
\definecolor{fmmistake}{HTML}{EF476F} % Common Mistake
\definecolor{fmnotice}{HTML}{9B5DE5}  % Notice
\definecolor{fmgray}{HTML}{6B6B6B}

\pagestyle{fancy}
\fancyhf{}
\renewcommand{\headrulewidth}{0pt}
\fancyfoot[C]{\footnotesize\color{fmgray}\textcopyright\ 2026 FangMath. All rights reserved.}
\setlength{\parindent}{0pt}

% ---- Card styles ----
\newtcolorbox{cardfront}[1]{enhanced, width=\linewidth, height=11.4cm,
  colback=white, colframe=fmtitle, boxrule=1.6pt, arc=4mm,
  valign=center, halign=center, left=10mm, right=10mm,
  overlay={\node[anchor=north west, font=\footnotesize\bfseries, text=fmtitle]
     at ([xshift=5mm,yshift=-4mm]frame.north west) {FLASH CARD 1};
     \node[anchor=south east, font=\scriptsize, text=fmgray]
     at ([xshift=-5mm,yshift=3mm]frame.south east) {Linear Functions \textperiodcentered\ Parallel \& Perpendicular Lines};}}
\newtcolorbox{cardback}[1]{enhanced, width=\linewidth, height=11.4cm,
  colback=white, colframe=fmkey, boxrule=1.6pt, arc=4mm,
  valign=top, left=6mm, right=6mm, top=4mm,
  overlay={\node[anchor=south east, font=\scriptsize, text=fmgray]
     at ([xshift=-5mm,yshift=3mm]frame.south east) {Card 1 \textperiodcentered\ Answer};}}

\newcommand{\cardtitle}[1]{{\large\bfseries\color{fmkey}#1}\par\vspace{1mm}
  {\color{fmkey}\rule{\linewidth}{0.8pt}}\par\vspace{2mm}}
\newcommand{\cutline}{\par\vspace{3mm}\noindent
  \begin{tikzpicture}\draw[dashed,fmgray] (0,0) -- (\linewidth-1pt,0);
  \node[fill=white,font=\scriptsize,text=fmgray] at (0.5\linewidth-0.5pt,0) {cut here};
  \end{tikzpicture}\par\vspace{3mm}}
\newcommand{\mistake}[1]{{\color{fmmistake}\bfseries Common Mistake:}\ #1}

% ---- Tiny pictures (two lines in a small window) ----
% \twolines{slope1}{int1}{slope2}{int2}
\newcommand{\twolines}[4]{%
  \begin{tikzpicture}[scale=0.4,baseline=(c.base)]
    \coordinate (c) at (0,-0.3);
    \clip (-2.2,-2.2) rectangle (2.2,2.2);
    \draw[gray!45,line width=0.4pt] (-2.2,0)--(2.2,0) (0,-2.2)--(0,2.2);
    \draw[fmtitle,line width=1.1pt,domain=-2.5:2.5] plot (\x,{#1*\x+#2});
    \draw[fmmistake,line width=1.1pt,domain=-2.5:2.5] plot (\x,{#3*\x+#4});
  \end{tikzpicture}}

\begin{document}
% =====================================================================
%  PAGE 1 -- FRONTS
% =====================================================================
\begin{cardfront}{1}
  \centering
  {\color{fmtitle}\Large\bfseries Slopes of Two Lines}\par\vspace{7mm}
  {\large Two lines lie in the \textbf{same plane}.\\[2mm]
  Their slopes are $m_1$ and $m_2$.}\par\vspace{7mm}
  {\Large\color{fmtitle} What are the \textbf{3 cases}?}\par\vspace{4mm}
  {\large\color{fmgray} What does the sign $\perp$ mean?}
\end{cardfront}

\cutline

\begin{tcolorbox}[enhanced, width=\linewidth, height=11.4cm,
  colback=white, colframe=fmtitle, boxrule=1.6pt, arc=4mm,
  valign=center, halign=center, left=10mm, right=10mm,
  overlay={\node[anchor=north west, font=\footnotesize\bfseries, text=fmtitle]
     at ([xshift=5mm,yshift=-4mm]frame.north west) {FLASH CARD 2};
     \node[anchor=south east, font=\scriptsize, text=fmgray]
     at ([xshift=-5mm,yshift=3mm]frame.south east) {Linear Functions \textperiodcentered\ Parallel \& Perpendicular Lines};}]
  \centering
  {\color{fmtitle}\Large\bfseries Writing a Parallel or\\[1mm] Perpendicular Line}\par\vspace{7mm}
  {\large Write the equation of the line through $(2,\,1)$ that is}\par\vspace{3mm}
  {\large \textbf{(a)} parallel to $y = 2x + 5$\qquad \textbf{(b)} perpendicular to $y = 2x + 5$.}\par\vspace{6mm}
  {\large\color{fmgray} Give your answers in slope-intercept form, $y = mx + b$.}
\end{tcolorbox}

\newpage
% =====================================================================
%  PAGE 2 -- BACKS (same top/bottom order as page 1)
% =====================================================================
\begin{cardback}{1}
\cardtitle{Slopes of Two Lines --- 3 Cases}
{\small Compare the slopes $m_1$ and $m_2$ of two lines in the same plane:}\par\vspace{2mm}
\renewcommand{\arraystretch}{1.6}
\begin{tabular}{@{}>{\centering\arraybackslash}m{2.0cm} m{3.4cm} m{3.4cm} m{5.4cm}@{}}
\textbf{\color{fmkey}Picture} & \textbf{\color{fmkey}Slopes} & \textbf{\color{fmkey}The lines are\ldots} & \textbf{\color{fmkey}Example}\\[1mm]
\twolines{1}{-0.9}{1}{0.9} & $m_1 = m_2$ & \textbf{parallel} & $y=3x-2$ and $y=3x+4$\newline {\footnotesize same slope $3$}\\
\twolines{2}{0}{-0.5}{0} & $m_1 \cdot m_2 = -1$\newline {\footnotesize opposite reciprocals} & \textbf{perpendicular} ($\perp$) & $y=2x+1$ and $y=-\tfrac12 x+4$\newline {\footnotesize $2 \cdot \left(-\tfrac12\right) = -1$}\\
\twolines{2}{0}{0.4}{0.4} & neither & \textbf{intersect, but\newline not perpendicular} & $y=2x$ and $y=x+3$\newline {\footnotesize $2 \neq 1$, and $2 \cdot 1 = 2 \neq -1$}
\end{tabular}\par\vspace{3mm}
{\small $\perp$ is the \textbf{perpendicular sign}: $\ell_1 \perp \ell_2$ means ``line $\ell_1$ is perpendicular to line $\ell_2$'' (they meet at a right angle).}\par\vspace{2mm}
{\footnotesize\color{fmnotice}\bfseries Watch out:}\ {\footnotesize
(1) Same slope \emph{and} same $y$-intercept means the \emph{same line}, not two parallel lines.\
(2) A vertical line ($x=a$) and a horizontal line ($y=b$) are perpendicular, even though a vertical line has no slope, so $m_1 \cdot m_2 = -1$ cannot be used.}
\end{cardback}

\cutline

\begin{tcolorbox}[enhanced, width=\linewidth, height=11.4cm,
  colback=white, colframe=fmkey, boxrule=1.6pt, arc=4mm,
  valign=top, left=6mm, right=6mm, top=4mm,
  overlay={\node[anchor=south east, font=\scriptsize, text=fmgray]
     at ([xshift=-5mm,yshift=3mm]frame.south east) {Card 2 \textperiodcentered\ Answer};}]
\cardtitle{Writing a Parallel or Perpendicular Line}
{\small\textbf{Steps:} \ \textbf{1.} Find the slope you need. \ \textbf{2.} Use point-slope form $y - y_1 = m(x - x_1)$. \ \textbf{3.} Solve for $y$.}\par\vspace{3mm}
\begin{minipage}[t]{0.48\linewidth}
{\color{fmwatch}\bfseries (a) Parallel}\par\vspace{1mm}
{\small Same slope as $y=2x+5$: \ $m = 2$.}
\begin{align*}
y - 1 &= 2(x - 2)\\
y - 1 &= 2x - 4\\
y &= 2x - 3
\end{align*}
\end{minipage}\hfill
\begin{minipage}[t]{0.48\linewidth}
{\color{fmwatch}\bfseries (b) Perpendicular}\par\vspace{1mm}
{\small Opposite reciprocal of $2$: \ $m = -\tfrac12$.}
\begin{align*}
y - 1 &= -\tfrac12(x - 2)\\
y - 1 &= -\tfrac12 x + 1\\
y &= -\tfrac12 x + 2
\end{align*}
\end{minipage}\par\vspace{2mm}
{\small\textbf{Check:} put in $x = 2$. \ (a) $2 \cdot 2 - 3 = 1$ \checkmark \qquad (b) $-\tfrac12 \cdot 2 + 2 = 1$ \checkmark}\par\vspace{3mm}
{\small\mistake{stopping at $y - 1 = 2x - 4$, or dropping the $-1$ and writing $y = 2x - 4$. Finish the last step: add $1$ to \emph{both} sides so $y$ is alone.}}
\end{tcolorbox}
\end{document}
